\chapter{Conclusions and Future Work}
\label{ch:conclusions}

This chapter draws together the findings of the thesis. The first section summarises the key quantitative results established in Chapter~3. The second section states the principal conclusions regarding Julia's viability as a high-performance alternative to Python for power system computation, addressing in turn each research question posed in the Introduction. The final section identifies a structured programme of future work that extends the implemented modules toward a broader Julia-based power system analysis framework.

% ──────────────────────────────────────────────────────────────────────────────
\section{Summary of Results}
% ──────────────────────────────────────────────────────────────────────────────

This thesis implemented and evaluated eight computational contributions: three core power system solvers ported from PyPSA to Julia (DC Power Flow, AC Power Flow, and single-period LOPF), an extension to multi-period LOPF with storage and wind, Unit Commitment (MILP), an LSTM load forecaster with conformal prediction intervals, a scenario-based stochastic LOPF, and a systematic validation campaign on standard IEEE test cases. The key findings are summarised below.

\textbf{Numerical correctness on 3-bus network.}
For the LP-based methods --- single-period LOPF, multi-period LOPF, locational marginal prices, and emission-constrained dispatch --- the Julia and PyPSA results are numerically identical (zero difference in generator dispatch, total cost, and nodal prices). For the power flow modules, active and reactive power flows agree to better than $10^{-3}$, while two quantities differ by a constant scale factor: the DC power flow voltage angles and the transformer branch flow. Both differences trace to a documented per-unit convention difference --- PyPSA normalises branch susceptance by $V_\text{nom}^2$ and the transformer reactance by its own rating $s_\text{nom}$, whereas the present implementation follows the MATPOWER convention of expressing susceptance in MW/rad on the system base. Crucially, these convention differences do not affect any economically meaningful quantity (dispatch, cost, or price), confirming the correctness of the porting.

\textbf{IEEE 14-bus standard test case validation.}
The implementations were additionally validated on the IEEE 14-bus benchmark network \cite{zimmerman2011matpower}. The AC power flow solution matches the MATPOWER reference to within $1.33 \times 10^{-3}$\,p.u.\ in voltage magnitude and $0.017^{\circ}$ in angle. The DC power flow and LOPF solutions from Julia and PyPSA are identical to numerical precision. This confirms that the implementations generalise correctly beyond the construction network.

\textbf{DC Power Flow performance.}
Measured as the noise-robust minimum with threads pinned, Julia's DC power flow speedup falls from $\sim$$6\times10^{4}$ on a 3-bus network to $169\times$ at 300 synthetic buses, and to $4.7\times$ on the 2{,}869-bus PEGASE network. The huge small-network ratios are PyPSA's fixed $\approx 120\,\text{ms}$ Python overhead, not a compute advantage; the honest compute-bound figure at scale is a low single-digit factor.

\textbf{AC Power Flow performance.}
AC power flow is the one method where Julia does \emph{not} hold a consistent advantage. The apparent speedup at small sizes ($24.6\times$ at 3 buses) is again PyPSA's fixed overhead; once amortised, the comparison reverses --- $1.5\times$ at 50 buses and $0.5\times$ at 100 buses, i.e.\ Julia is about twice as slow. The cause is algorithmic: Julia delegates AC power flow to the Ipopt interior-point solver (PowerModels.jl), whose per-iteration cost exceeds PyPSA's specialised Newton--Raphson, and the model-assembly advantage that drives the LP results has nothing to act on when a third-party nonlinear solver dominates.

\textbf{Single-period LOPF performance.}
The Julia LOPF (JuMP + HiGHS) is $118$--$283\times$ faster on synthetic networks and $58$--$117\times$ faster on the standard IEEE/PEGASE cases; on the 2{,}869-bus PEGASE grid PyPSA takes $28.4$\,s against Julia's $0.49$\,s ($58\times$). Since both implementations use the same HiGHS LP solver, the entire speedup comes from JuMP's compiled model assembly versus PyPSA/linopy's Python constraint construction. The speedup declines gently with size but remains large at scale, which is directly applicable to large-scale studies such as PyPSA-Eur.

\textbf{Multi-period LOPF with storage and wind.}
The 24-hour multi-period LOPF with a StorageUnit and variable wind generation was successfully implemented and validated. Both Julia and PyPSA find the same optimal cost of $126{,}790.79\,\text{\euro}$. The LP correctly exploits time-arbitrage: the storage charges at night ($\approx 22\,\text{\euro/MWh}$ round-trip) and discharges at peak to displace expensive backup generation ($50\,\text{\euro/MWh}$). Performance benchmarks over horizons $T \in \{6, 12, 24, 48\}$ hours show Julia achieving $89$--$298\times$ speedups. The PyPSA solve time is nearly constant ($\approx 1.14\,\text{s}$) across all horizons, a direct demonstration that PyPSA's bottleneck is model assembly rather than solving; Julia's solve time scales linearly with $T$, from $3.8$\,ms at $T = 6$ to $12.8$\,ms at $T = 48$.

\textbf{Unit Commitment.}
The MILP Unit Commitment solver was implemented and benchmarked. Julia achieves $17$--$79\times$ speedups over PyPSA across horizons $T \in \{6, 12, 24\}$ and network sizes of 3--30 buses. The speedup is lower than for continuous LOPF because the MILP branch-and-bound search --- executed by the same HiGHS solver in both implementations --- constitutes the dominant cost. LMPs are extracted via LP relaxation with binaries fixed, following the standard electricity market methodology~\cite{oneill2005}.

\textbf{AI component: LSTM forecasting and stochastic dispatch.}
An LSTM load forecaster (Flux.jl, hidden size 64, weekly 168\,h window, 24\,h horizon, residual-to-seasonal-naive target with conformal calibration) was trained and evaluated on \emph{real} German load (OPSD/ENTSO-E) with a strict 60-day out-of-time hold-out, averaged over three seeds. The honest finding is reported as it stands: the LSTM beats the persistence baseline ($+42.7\%$ RMSE skill, $+61.5\%$ with temperature) but \emph{does not} beat the seasonal-naive baseline (MAPE $5.83\%$ with calendar features, $3.67\%$ with idealized temperature, versus $2.77\%$ for seasonal-naive). For aggregate national load the weekly cycle is nearly deterministic, so a learned model adds value only through an exogenous signal; temperature is confirmed to be that decisive lever, closing the gap from $-84\%$ to $-23\%$ RMSE skill.

The contribution of the AI component is therefore not point-forecast accuracy but the \emph{uncertainty quantification} it enables. The conformal intervals are propagated into a scenario-based stochastic LOPF (Sample Average Approximation), which on the demonstration network returns an expected cost, a worst-case CVaR$_{90}$, and a cost spread ($\sigma$). The gap between expected and worst-case cost ($\approx 13\%$) is the reserve a risk-averse operator must budget for the worst $10\%$ of demand outcomes --- a decision-relevant risk measure that no deterministic dispatch, however accurate its point forecast, can provide.

This AI component is presented as a methodological contribution rather than a production forecaster. Its honest limitation --- a neural forecaster that loses to a trivial seasonal baseline on point accuracy --- is itself an informative result: it locates the value of machine learning in this domain in calibrated uncertainty and weather-informed features rather than in architecture alone, and it motivates the future work in Section~\ref{sec:future}.

% ──────────────────────────────────────────────────────────────────────────────
\section{Conclusions}
% ──────────────────────────────────────────────────────────────────────────────

The experimental results demonstrate that porting PyPSA's core algorithms to Julia is both feasible and highly beneficial. The main conclusions are:

\begin{enumerate}
    \item \textbf{Julia is a viable replacement for Python in power system computation.}
    The mathematical formulations translate directly from PyPSA's Python code, and the Julia ecosystem (JuMP, HiGHS, PowerModels.jl) provides mature, well-maintained tools for every required component. Validation on both a custom 3-bus network and the standard IEEE 14-bus case confirms numerical equivalence.

    \item \textbf{The performance advantage is largest for LP-based optimisation.}
    The LP-based methods (single- and multi-period LOPF, unit commitment) are tens to a few hundred times faster, and the advantage persists at scale ($58\times$ for LOPF on the 2{,}869-bus PEGASE grid, where PyPSA takes $28$\,s and Julia under $0.5$\,s). Because both sides call the same HiGHS solver, this gain is entirely the compiled model-assembly layer --- making Julia particularly valuable for iterative large-scale studies such as PyPSA-Eur.

    \item \textbf{Multi-period optimisation with storage scales naturally.}
    Extending the single-period LOPF to 24 time periods with a StorageUnit required only incremental additions to the JuMP model --- SOC dynamics, charge/discharge variables, and a cyclicity constraint. The Julia implementation produces results identical to PyPSA, confirming that the time-coupling constraints are correctly formulated.

    \item \textbf{Algorithm choice matters more than language for AC power flow.}
    AC power flow is the honest exception: Julia is competitive only at small sizes and becomes about twice as slow as PyPSA at 100 buses, because it delegates to the Ipopt interior-point solver while PyPSA uses a specialised Newton--Raphson iteration. The compiled model-assembly advantage that drives the LP speedups does not apply when a third-party nonlinear solver dominates the work. A native Newton--Raphson implementation in Julia would be required to recover an advantage here.

    \item \textbf{JIT warmup is not a practical barrier.}
    For realistic use cases --- time-series analysis, scenario sweeps, iterative optimisation --- functions are invoked thousands of times, making the one-time JIT compilation cost ($2$--$5\,\text{s}$) negligible. The benchmarks in this thesis measure steady-state (post-warmup) performance, which is the relevant metric for production use.

    \item \textbf{The existing Julia power system landscape has a gap.}
    As documented in the literature review (Chapter~\ref{ch:litreview}), no actively maintained, self-contained Julia reimplementation of PyPSA exists. PSA.jl and EnergyModels.jl were abandoned years ago; PowerModels.jl and PowerSimulations.jl follow independent architectures. This thesis provides a validated starting point for filling that gap.

    \item \textbf{AI-assisted dispatch quantifies risk, even when the forecaster does not win on accuracy.}
    On real load data the LSTM does not beat the seasonal-naive baseline on point accuracy --- a result reported honestly. Its value lies elsewhere: the conformal prediction framework provides distribution-free $90\%$-coverage intervals with no assumption on the residual distribution, and propagating these into the stochastic SAA gives operators a CVaR$_{90}$ tail-risk estimate that is inaccessible to any deterministic formulation, however accurate its point forecast. The AI component thus demonstrates that the value of machine learning in this pipeline is calibrated uncertainty quantification feeding risk-aware dispatch, not raw forecast accuracy.
\end{enumerate}

% ──────────────────────────────────────────────────────────────────────────────
\section{Future Work}
\label{sec:future}
% ──────────────────────────────────────────────────────────────────────────────

This thesis provides a foundation for a broader Julia-based power system analysis framework. Several concrete directions for future work are identified.

\textbf{Native Newton--Raphson AC power flow.}
The current AC power flow implementation delegates to Ipopt, which makes Julia slower than PyPSA's specialised Newton--Raphson at 100 buses. Implementing a direct Newton--Raphson solver in Julia --- building and factorising the Jacobian matrix natively --- would remove this algorithmic penalty and is the most impactful single improvement, as it is the only method where the present implementation does not already outperform PyPSA.

\textbf{Parallel stochastic LOPF.}
The SAA implemented in Section~\ref{sec:ai_component} solves $S$ independent LOPFs sequentially. Each scenario solve is fully independent, making the SAA embarrassingly parallel. Exploiting \texttt{Threads.@threads} or \texttt{Distributed.pmap} in Julia would reduce stochastic LOPF wall-clock time by a factor of $S$ on a multi-core machine, enabling real-time stochastic dispatch.

\textbf{Numerical validation on larger standard cases.}
Performance was already benchmarked on the IEEE 118/300 and PEGASE 1354/2869 networks, and numerical validation against PyPSA was performed on the 3-bus and IEEE 14-bus cases. A natural extension is a full variable-by-variable numerical validation against PyPSA on the larger PEGASE cases, and eventually on a subset of the PyPSA-Eur European transmission network \cite{horsch2018pypsa_eur}, to extend the equivalence guarantee to real-world problem sizes.

\textbf{Weather-informed and residual forecasting.}
The evaluation showed that the LSTM matches the seasonal-naive baseline only with idealized (perfect-foresight) temperature. Coupling the forecaster to an actual numerical weather prediction feed, and combining it with the strong seasonal-naive baseline in an explicit residual model, is the most promising route to a forecaster that genuinely improves on the naive baselines on real out-of-time data.

\textbf{Security-constrained OPF.}
Adding N-1 contingency constraints for transmission line or generator outages would produce a security-constrained LOPF (SCLOPF), which is a standard requirement for transmission system operator studies.

\textbf{Capacity expansion planning.}
Extending the multi-period LOPF with investment variables (\texttt{p\_nom\_extendable}) would enable greenfield or retrofit generation planning studies, equivalent to PyPSA's capacity expansion mode.

\textbf{PyPSA data format interoperability.}
Developing a Julia reader for PyPSA's native netCDF/HDF5 network format would enable seamless use of existing PyPSA-Eur model files in the Julia solvers, without requiring manual data conversion.

The results of this thesis demonstrate that a Julia-based power system analysis framework can deliver order-of-magnitude performance improvements over Python while maintaining numerical equivalence with PyPSA, and that integrating a data-driven forecasting layer enables uncertainty-aware, risk-quantified dispatch decisions that are inaccessible to deterministic formulations.
